% Generated draft: requires mathematical and source review.
\documentclass[11pt]{article}
\usepackage[margin=25mm]{geometry}
\usepackage{fontspec}
\setmainfont{Latin Modern Roman}
\usepackage{amsmath,amssymb,mathtools,unicode-math}
\setmathfont{Latin Modern Math}
\usepackage{enumitem,needspace,xurl,hyperref}
\setlength{\parindent}{0pt}
\setlength{\parskip}{5pt}
\setlength{\emergencystretch}{3em}
\newcommand{\N}{\mathbb{N}}
\newcommand{\Z}{\mathbb{Z}}
\newcommand{\Q}{\mathbb{Q}}
\newcommand{\R}{\mathbb{R}}
\newcommand{\C}{\mathbb{C}}
\newcommand{\F}{\mathbb{F}}
\newcommand{\A}{\mathbb{A}}
\newcommand{\PP}{\mathbb{P}}
\newcommand{\E}{\mathbb{E}}
\newcommand{\eps}{\varepsilon}
\newcommand{\dd}{\,\mathrm{d}}
\begin{document}
\section*{UnsolvedMath ID: 30002572}
\textbf{Source problem code:} catalogue entry
\subsection{Definitions and mathematical statement}
The supplied record describes a collection of research questions about cancer-cell migration, rather than a single specified mathematical conjecture. Its terms identify modeling topics: microscopic and mesoscopic models describe migration at different levels of resolution; a macroscopic limit is a limiting description at a larger scale; and an inverse recovery problem seeks an unknown model quantity from observations. These descriptions do not specify equations, admissible models, observations, or limiting procedures.

The research tasks recorded under permanent UnsolvedMath ID \(30002572\), source code \texttt{catalogue entry}, are:
\begin{enumerate}
\item Describe physical cell passage through obstacles.
\item Determine the mechanics of collective motion and develop models of that motion.
\item Optimize the methods identified in the source as DTI and Q-ball methods.
\item Establish benchmarks for movement models.
\item Formulate and analyze the moving-boundary problem for acid-mediated invasion.
\item Compare amoeboid, mesenchymal, and collective migration strategies, including their invasion speeds.
\item Establish global bounds for chemotaxis--haptotaxis models with tissue remodeling.
\item Recover chemotaxis sensitivity from observations through an inverse problem.
\item Develop simple microscopic or mesoscopic models of collective motion, including leaderless motion, and determine their macroscopic limits.
\end{enumerate}

These tasks retain the scope of the supplied statement, but they do not constitute fully specified mathematical assertions. In particular, the record supplies no obstacle geometry or cell mechanics; collective-motion interaction rules; imaging optimization criterion; benchmark data or comparison criterion; acid-mediated invasion equations or interface conditions; definition of invasion speed or common conditions for comparing migration strategies; chemotaxis--haptotaxis equations, parameter hypotheses, or norms for global bounds; observation operator or admissible sensitivity class for inverse recovery; or scaling and convergence notion for macroscopic limits. Spatial domains, initial and boundary conditions, solution concepts, and quantifiers are also unspecified. Consequently, no particular existence, boundedness, uniqueness, identifiability, optimization, or convergence theorem is asserted here. Choosing those missing specifications requires source clarification and human mathematical review.
\subsection{Short English statement}
This record collects nine research directions in cancer-cell migration, spanning mechanics, imaging, model comparison, invasion, global bounds, inverse recovery, and collective-motion limits. It lacks the equations and hypotheses needed for precise theorem statements. Its open-status label and partial-results assessment are supplied dated curation, not independently verified conclusions.
\subsection{Sources}
\begin{itemize}
\item Local UnsolvedMath registry: \texttt{data/unsolvedmath/problems.json}, permanent ID \(30002572\), title ``Multiscale Problems in Cancer-Cell Migration,'' source code \texttt{catalogue entry}, collection ``Oberwolfach Reports Open Problems.'' The supplied record gives proposed year \(2014\), no named proposer, and status \texttt{open}.
\item Supplied report attribution: \emph{Mini-Workshop: Mathematical Models for Cancer Cell Migration} (2015), \emph{Oberwolfach Reports}, pp.\ 1075--1109, DOI \texttt{10.4171/owr/2014/19}. \url{https://doi.org/10.4171/owr/2014/19}. The background describes Oberwolfach Reports as documenting mathematical workshops and their open-problem discussions. No individual problem author or problem-specific page is supplied.
\item Supplied literature source, labeled ``Migration-model comparison review'': \url{https://doi.org/10.1371/journal.pcbi.1008411}.
\item Supplied literature source, labeled ``Chemotaxis--haptotaxis global existence'': \url{https://arxiv.org/abs/1711.10044}.
\item Literature assessment dated exactly \texttt{2026-08-21}: the supplied curation labels the problem open and reports partial results for individual subquestions, including boundedness and global-existence theorems for restricted chemotaxis--haptotaxis systems and inverse identification methods. It states that the listed benchmark, imaging, invasion-speed, and general-bound questions remain unresolved as a package. This assessment is supplied dated curation derived from the supplied OWR curation, not independently verified; time-sensitive status and cited claims require independent verification.
\end{itemize}
\end{document}
